% ================================================
% PRÉAMBULE COMPLET - COURS POO JAVA
% ================================================

\documentclass[12pt,a4paper]{article}

% ================================================
% PACKAGES ESSENTIELS
% ================================================

% Langue et encodage
\usepackage[utf8]{inputenc}
\usepackage[T1]{fontenc}
\usepackage[french]{babel}

% Mise en page
\usepackage[margin=2.5cm]{geometry}
\usepackage{fancyhdr}
\usepackage{lastpage}

% Mathématiques
\usepackage{amsmath, amssymb, amsthm}

% Graphiques et couleurs
\usepackage{graphicx}
\usepackage{xcolor}
\usepackage{tikz}

% Tableaux
\usepackage{array}
\usepackage{booktabs}
\usepackage{multirow}

% Listes
\usepackage{enumitem}

% Code source
\usepackage{listings}
\usepackage{minted} % Alternative (nécessite Python + Pygments)

% Boîtes colorées
\usepackage{tcolorbox}
\tcbuselibrary{listings, skins, breakable}

% Liens hypertextes
\usepackage{hyperref}

% Icônes
\usepackage{fontawesome5}

% QR Codes (optionnel)
\usepackage{qrcode}

% ================================================
% CONFIGURATION HYPERREF
% ================================================

\hypersetup{
    colorlinks=true,
    linkcolor=blue,
    filecolor=magenta,      
    urlcolor=cyan,
    pdftitle={Cours POO Java - Programmation Orientée Objet},
    pdfauthor={Votre Nom},
    pdfsubject={Programmation Orientée Objet en Java},
    pdfkeywords={Java, POO, Cours, Programmation},
    pdfcreator={LaTeX}
}

% ================================================
% COULEURS PERSONNALISÉES
% ================================================

\definecolor{githubblue}{RGB}{0, 119, 200}
\definecolor{githubgray}{RGB}{36, 41, 46}
\definecolor{javagreen}{RGB}{0, 128, 0}
\definecolor{javapurple}{RGB}{128, 0, 128}
\definecolor{javablue}{RGB}{0, 0, 255}
\definecolor{javared}{RGB}{255, 0, 0}
\definecolor{codebackground}{RGB}{248, 248, 248}

% ================================================
% CONFIGURATION LISTINGS (CODE JAVA)
% ================================================

\lstdefinestyle{javastyle}{
    language=Java,
    backgroundcolor=\color{codebackground},
    basicstyle=\ttfamily\small,
    keywordstyle=\color{javablue}\bfseries,
    stringstyle=\color{javared},
    commentstyle=\color{javagreen}\itshape,
    numberstyle=\tiny\color{gray},
    numbers=left,
    numbersep=8pt,
    tabsize=4,
    showspaces=false,
    showstringspaces=false,
    showtabs=false,
    breaklines=true,
    breakatwhitespace=true,
    frame=single,
    rulecolor=\color{gray},
    captionpos=b,
    xleftmargin=15pt,
    xrightmargin=5pt,
    aboveskip=10pt,
    belowskip=10pt
}

\lstset{style=javastyle}

% ================================================
% BOÎTES TCOLORBOX PERSONNALISÉES
% ================================================

% Boîte GitHub
\newtcolorbox{githubbox}[1]{
    colback=githubblue!5!white,
    colframe=githubblue,
    fonttitle=\bfseries,
    title={\faGithub\ #1},
    boxrule=0.8pt,
    arc=3mm,
    left=8pt,
    right=8pt,
    top=8pt,
    bottom=8pt,
    breakable
}

% Boîte Définition
\newtcolorbox{definition}[1]{
    colback=green!5!white,
    colframe=green!75!black,
    fonttitle=\bfseries,
    title={\faBook\ Définition : #1},
    boxrule=0.8pt,
    arc=3mm,
    breakable
}

% Boîte Remarque
\newtcolorbox{remarque}{
    colback=orange!5!white,
    colframe=orange!75!black,
    fonttitle=\bfseries,
    title={\faExclamationTriangle\ Remarque},
    boxrule=0.8pt,
    arc=3mm,
    breakable
}

% Boîte Exemple
\newtcolorbox{exemple}[1]{
    colback=blue!5!white,
    colframe=blue!75!black,
    fonttitle=\bfseries,
    title={\faLightbulb\ Exemple : #1},
    boxrule=0.8pt,
    arc=3mm,
    breakable
}

% Boîte Exercice
\newtcolorbox{exercice}[1]{
    colback=purple!5!white,
    colframe=purple!75!black,
    fonttitle=\bfseries,
    title={\faPencilAlt\ Exercice #1},
    boxrule=0.8pt,
    arc=3mm,
    breakable
}

% ================================================
% COMMANDES PERSONNALISÉES GITHUB
% ================================================

% URL du dépôt
\newcommand{\repourl}{https://github.com/GUEDDES/Cours-POO-Java}

% Lien GitHub simple
\newcommand{\githublink}[1]{%
    \href{\repourl/#1}{\texttt{#1}}%
}

% Lien GitHub avec icône
\newcommand{\githubrepo}{%
    \faGithub\ \href{\repourl}{github.com/GUEDDES/Cours-POO-Java}%
}

% Commande pour encart "Voir GitHub"
\newcommand{\voirgithub}[2]{%
    \begin{githubbox}{Ressources GitHub}
        \textbf{#1}
        
        \vspace{3pt}
        
        Voir sur GitHub : \href{\repourl/#2}{\texttt{#2}}
    \end{githubbox}
}

% Commande pour lien vers exemple
\newcommand{\githubexemple}[2]{%
    \href{\repourl/blob/main/#1}{\faCode\ #2}%
}

% Commande pour lien vers exercice
\newcommand{\githubexercice}[2]{%
    \href{\repourl/blob/main/#1}{\faPencilAlt\ #2}%
}

% ================================================
% EN-TÊTE ET PIED DE PAGE
% ================================================

\pagestyle{fancy}
\fancyhf{} % Clear default
\fancyhead[L]{\textbf{Cours POO Java}}
\fancyhead[R]{\leftmark}
\fancyfoot[L]{\githubrepo}
\fancyfoot[C]{Page \thepage\ sur \pageref{LastPage}}
\fancyfoot[R]{\today}
\renewcommand{\headrulewidth}{0.4pt}
\renewcommand{\footrulewidth}{0.4pt}

% ================================================
% COMMANDES UTILES
% ================================================

% Commande Java inline
\newcommand{\java}[1]{\texttt{#1}}

% Classe Java
\newcommand{\javaclass}[1]{\texttt{\textbf{#1}}}

% Méthode Java
\newcommand{\javamethod}[1]{\texttt{#1()}}

% ================================================
% TITRE DU DOCUMENT
% ================================================

\title{%
    \huge\textbf{Programmation Orientée Objet} \\[0.5cm]
    \Large Cours Complet avec Java \\[0.3cm]
    \large Année Universitaire 2024-2025
}

\author{Votre Nom}
\date{\today}

% ================================================
% FIN DU PRÉAMBULE
% ================================================
